\section{Cross-Level Analysis}
\label{sec:cross_level}
\label{sec:module-declaration}

Sections~\ref{sec:module-import}--\ref{sec:namespace-graph} analyzed the library at three distinct granularities: source files (\S\ref{sec:module-import}), declarations (\S\ref{sec:theorem-premise}), and namespaces (\S\ref{sec:namespace-graph}). Each level reveals a different facet of the tension between human organization and mathematical logic. We now overlay these levels, measuring how the three organizational structures---file boundaries, namespace boundaries, and logical dependencies---align and diverge.

The overlay yields three cross-level measurements. Namespace--module alignment (\S\ref{sec:ns-mod-cross}) measures the agreement between the two human organizational systems---naming and filing. Module cohesion (\S\ref{sec:cohesion}) quantifies how much declaration-level dependency traffic stays within file boundaries. Namespace--declaration interaction (\S\ref{sec:ns-decl-interaction}) compares the constraining power of the two systems at the declaration level and reveals a directional depth asymmetry between the import and declaration graphs. Assembling these measurements (\S\ref{sec:three-layers}) reveals a consistent pattern: human organizational decisions are internally coherent but collectively divergent from mathematical logic.

\subsection{Namespace--Module Alignment}
\label{sec:ns-mod-cross}

The containment analysis of \S\ref{sec:containment-decay} and the cohesion analysis below each compare one human organizational system against the logical dependency structure. We first ask a different question: how well do the two \emph{human} systems---naming and filing---agree with each other? If both reflect the same cognitive model, their alignment should be high even when both diverge from logic.

Table~\ref{tab:ns-mod-cross} cross-tabulates the namespace and module partitions for the $217{,}487$ declarations with file mappings. The normalized mutual information (NMI) between the two partitions is $0.708$ at depth~$1$ and $0.766$ at depth~$2$---substantially higher than the NMI of~$0.34$ between Louvain communities and the namespace hierarchy (\S\ref{sec:thm-community}).

\begin{table}[ht]
\centering
\caption{Namespace--module cross-tabulation for $217{,}487$ declarations with file mappings.}
\label{tab:ns-mod-cross}
\begin{tabular}{lrr}
\toprule
Metric & Depth $1$ & Depth $2$ \\
\midrule
Unique namespaces              & $2{,}241$  & $4{,}943$ \\
Unique modules (files)         & \multicolumn{2}{c}{$6{,}993$} \\
NMI(namespace, module)         & $0.708$    & $0.766$ \\
Single-file namespaces         & $50\%$     & $64\%$ \\
Single-namespace files         & $58\%$     & $39\%$ \\
Max files per namespace        & \multicolumn{2}{c}{$2{,}169$ (\ns{\_root\_})} \\
Max namespaces per file        & $31$       & $34$ \\
\bottomrule
\end{tabular}
\end{table}

At depth~$1$, half of all namespaces appear in exactly one file---for these, naming and filing are perfectly aligned. But the tail is heavy: \ns{CategoryTheory} spans $977$ files, and~$10\%$ of namespaces span more than $10$ files (Table~\ref{tab:ns-mod-cross}). In the reverse direction, $58\%$ of files contain declarations from only one namespace; the most namespace-diverse files are concentrated in topology: \module{Topology.Constructions} and \module{Topology.Separation.Basic} each host over $30$ depth-$1$ namespaces, reflecting their role as bridge files that simultaneously manipulate declarations from many conceptual domains.

The two directions of misalignment reflect distinct engineering pressures. When a namespace spans many files (namespace~$>$~module), a single conceptual domain has been partitioned into multiple compilation units---typically to manage build times. When a file contains many namespaces (module~$>$~namespace), related concepts from different naming domains have been gathered in one place for cognitive convenience. Both patterns are rational organizational decisions; they produce NMI~$\approx 0.7$ rather than~$1.0$ because the two pressures---compilation efficiency and cognitive proximity---do not always agree.

\begin{remark}
\label{rem:two-pressures}
The NMI of~$0.71$ between namespaces and modules is \emph{twice} the NMI of~$0.34$ between namespaces and dependency-based Louvain communities (\S\ref{sec:thm-community}). Human organizational decisions---naming and filing---agree with each other far more than either agrees with the logical structure of the dependency graph. This is unsurprising: both naming and filing are human decisions, often made by the same contributor, while the dependency structure is determined by mathematical necessity.
\end{remark}

\subsection{Module--Declaration Interaction}
\label{sec:cohesion}

In $G_{\mathrm{thm}}$, each module becomes a \emph{region}---a set of declaration-vertices sharing a common source file. For each module~$m$, let $\mathcal{D}_m \subseteq \mathcal{D}$ denote its declarations. We partition the incident edges into two classes:
\begin{align*}
  E_{\mathrm{int}}(m) &= \{(d_1, d_2) \in E_{\mathrm{thm}} \mid d_1, d_2 \in \mathcal{D}_m\}, \\
  E_{\mathrm{ext}}(m) &= \{(d_1, d_2) \in E_{\mathrm{thm}} \mid d_1 \in \mathcal{D}_m \oplus d_2 \in \mathcal{D}_m\},
\end{align*}
where $\oplus$ denotes exclusive or. Figure~\ref{fig:module-region} illustrates this partition: of six dependency edges among declarations in \module{Data.Nat.Basic}, \module{Data.Nat.Defs}, and \module{Init.Prelude}, only one stays within its module---the rest pierce through file boundaries.

\begin{figure}[H]
\centering
\begin{subfigure}[t]{0.36\textwidth}
\centering
\vspace{0pt}
% ── Directory structure (same style as Figure 1(a)) ──
\begin{forest}
  for tree={
    font=\scriptsize\ttfamily,
    text=blue!60!black,
    grow'=0,
    child anchor=west,
    parent anchor=south,
    anchor=west,
    calign=first,
    edge path={
      \noexpand\path [draw, black, \forestoption{edge}]
        (!u.south west) +(7.5pt,0) |- (.child anchor)\forestoption{edge label};
    },
    before typesetting nodes={
      if n=1 {insert before={[,phantom]}} {}
    },
    fit=band,
    before computing xy={l=12pt},
  }
[Data/
  [Nat/
    [Basic.lean
      [\textcolor{red!70!black}{add\_comm}]
      [\textcolor{red!70!black}{add\_left\_comm}]
    ]
    [Defs.lean
      [\textcolor{red!70!black}{succ\_eq\_add\_one}]
      [\textcolor{red!70!black}{zero\_add}]
    ]
  ]
]
\end{forest}
\caption{File-system hierarchy~$T$ (\textcolor{blue!60!black}{blue}), extended to the declaration level.}
\label{fig:tree-view}
\end{subfigure}%
\hfill
% ── Panel (b): Source code (top right, next to tree) ──
\begin{subfigure}[t]{0.60\textwidth}
\centering
\vspace{0pt}
\raggedright\scriptsize
\textsf{\bfseries\color{blue!60!black}Basic.lean}\\[2pt]
\ttfamily
\textbf{import} \textcolor{blue!60!black}{Init.Prelude}\\[3pt]
\textbf{theorem} \textcolor{red!70!black}{add\_comm} (n m : Nat)\\
\hspace{1.5em}: n + m = m + n := \textbf{by}\\
\hspace{1.5em}induction n <;> simp [\textcolor{red!70!black}{Eq.refl}]\\[3pt]
\textbf{theorem} \textcolor{red!70!black}{add\_left\_comm} (n m k : Nat)\\
\hspace{1.5em}: n + (m + k) = m + (n + k) := \textbf{by}\\
\hspace{1.5em}rw [\textcolor{red!70!black}{add\_comm}]; simp [\textcolor{red!70!black}{Eq.refl}]\\[6pt]
\textsf{\bfseries\color{blue!60!black}Defs.lean}\\[2pt]
\ttfamily
\textbf{import} \textcolor{blue!60!black}{Data.Nat.Basic}\\
\textbf{import} \textcolor{blue!60!black}{Init.Prelude}\\[3pt]
\textbf{theorem} \textcolor{red!70!black}{succ\_eq\_add\_one} (n : Nat)\\
\hspace{1.5em}: succ n = n + 1 := \textbf{by}\\
\hspace{1.5em}rw [\textcolor{red!70!black}{add\_comm}]; exact \textcolor{red!70!black}{Eq.refl} \_\\[3pt]
\textbf{theorem} \textcolor{red!70!black}{zero\_add} (n : Nat)\\
\hspace{1.5em}: 0 + n = n := \textbf{by}\\
\hspace{1.5em}induction n <;> exact \textcolor{red!70!black}{Eq.refl} \_
\caption{Proof source code: each \textcolor{red!70!black}{red identifier} in a proof term generates a dependency edge in~(c).}
\label{fig:source-view}
\end{subfigure}

\par\vspace{4pt}

% ── Panel (c): Isometric dependency graph (full width, bottom) ──
\begin{subfigure}[b]{\textwidth}
\centering
\begin{tikzpicture}[
  every node/.style={font=\scriptsize\ttfamily, inner sep=2pt},
  dot/.style={circle, fill=red!70!black, minimum size=3.5pt, inner sep=0pt},
  dep/.style={->, >=Stealth, red!70!black, semithick},
  depext/.style={->, >=Stealth, red!70!black, semithick, dashed},
  modimp/.style={->, >=Stealth, blue!60!black, thick},
]
  % ── Card 3 (top, back): module Defs ──
  \filldraw[fill=blue!8, draw=blue!50!black, semithick, line join=round]
    (0.2,4.7) -- (10.2,4.7) -- (11.4,5.09) -- (1.4,5.09) -- cycle;
  \filldraw[fill=blue!18, draw=blue!50!black, semithick, line join=round]
    (0.2,4.7) -- (10.2,4.7) -- (10.2,4.57) -- (0.2,4.57) -- cycle;
  \node[font=\tiny\sffamily, blue!60!black, anchor=west] at (10.4,4.63)
    {\textnormal{\module{Data.Nat.Defs}}};
  \node[dot] (sea) at (3.5,4.88) {};
  \node[font=\tiny\ttfamily, red!70!black, anchor=west] at (3.6,4.88) {succ\_eq\_add\_one};
  \node[dot] (za)  at (7.5,4.88) {};
  \node[font=\tiny\ttfamily, red!70!black, anchor=west] at (7.6,4.88) {zero\_add};

  % ── Card 2 (middle): module Basic ──
  \filldraw[fill=blue!8, draw=blue!50!black, semithick, line join=round]
    (0,2.5) -- (10.0,2.5) -- (11.2,2.89) -- (1.2,2.89) -- cycle;
  \filldraw[fill=blue!18, draw=blue!50!black, semithick, line join=round]
    (0,2.5) -- (10.0,2.5) -- (10.0,2.37) -- (0,2.37) -- cycle;
  \node[font=\tiny\sffamily, blue!60!black, anchor=west] at (10.2,2.43)
    {\textnormal{\module{Data.Nat.Basic}}};
  \node[dot] (ac)  at (3.3,2.68) {};
  \node[font=\tiny\ttfamily, red!70!black, anchor=east] at (3.2,2.68) {add\_comm};
  \node[dot] (alc) at (7.3,2.68) {};
  \node[font=\tiny\ttfamily, red!70!black, anchor=west] at (7.4,2.68) {add\_left\_comm};

  % ── Card 1 (bottom, front): module Init.Prelude ──
  \filldraw[fill=blue!8, draw=blue!50!black, semithick, line join=round]
    (1.0,0.15) -- (8.0,0.15) -- (8.9,0.46) -- (1.9,0.46) -- cycle;
  \filldraw[fill=blue!18, draw=blue!50!black, semithick, line join=round]
    (1.0,0.15) -- (8.0,0.15) -- (8.0,0.03) -- (1.0,0.03) -- cycle;
  \node[font=\tiny\sffamily, blue!60!black, anchor=west] at (8.2,0.09)
    {\textnormal{\module{Init.Prelude}}};
  \node[dot] (er)  at (4.8,0.3) {};
  \node[font=\tiny\ttfamily, red!70!black, anchor=west] at (4.9,0.3) {Eq.refl};

  % ── Namespace bracket (Data.Nat spans Cards 2–3) ──
  \draw[blue!50, semithick]
    (-0.2,2.37) -- (-0.5,2.37) -- (-0.5,5.09) -- (-0.2,5.09);
  \node[font=\tiny\sffamily, blue!60!black, rotate=90, anchor=south]
    at (-0.75,3.73) {namespace \textnormal{\ns{Data.Nat}}};

  % ── Module import dots & arrows (blue = $G_{\mathrm{module}}$) ──
  \node[circle, fill=blue!60!black, minimum size=5.5pt, inner sep=0pt] (md) at (9.8, 4.88) {};
  \node[circle, fill=blue!60!black, minimum size=5.5pt, inner sep=0pt] (mb) at (10.0, 2.68) {};
  \node[circle, fill=blue!60!black, minimum size=5.5pt, inner sep=0pt] (mo) at (7.2, 0.3) {};
  \draw[modimp] (md) to[out=-85,in=85] (mb);                         % Defs → Basic
  \draw[modimp] (mb) to[out=-70,in=40] (mo);                         % Basic → other ns
  \draw[modimp] (md) to[out=-50,in=55] (mo);                         % Defs → other ns

  % ── Dependency arrows (red = mathematical logic) ──
  % Intra-module: the ONLY arrow that stays on its card
  \draw[dep] (alc) -- (ac);
  % Cross-module: Card 3 → Card 2 (external)
  \draw[depext] (sea) to[out=-88,in=92] (ac);
  % Cross-namespace: all → Eq.refl (external, piercing through card boundaries)
  \draw[depext] (ac)  to[out=-88,in=105] (er);
  \draw[depext] (alc) to[out=-88,in=60] (er);
  \draw[depext] (sea) to[out=-85,in=120] (er);
  \draw[depext] (za)  to[out=-85,in=45] (er);

\end{tikzpicture}
\caption{An isometric view of $G_{\mathrm{thm}}$ layered by module. Each \textcolor{blue!60!black}{blue card} represents a module file; \textcolor{red!70!black}{red nodes} are declarations. Solid \textcolor{red!70!black}{red arrows} show intra-module edges ($E_{\mathrm{int}}$); dashed \textcolor{red!70!black}{red arrows} show external edges ($E_{\mathrm{ext}}$) that pierce module boundaries. \textcolor{blue!60!black}{Blue arrows} on the right show the corresponding module-level imports from $G_{\mathrm{module}}$.}
\label{fig:dep-view}
\end{subfigure}
\caption{Three views of the same declarations. (a)~The file-system hierarchy~$T$ (\textcolor{blue!60!black}{blue}), extended to the declaration level. (b)~Proof source code: each \textcolor{red!70!black}{red identifier} creates a dependency edge in~(c). (c)~An isometric view of~$G_{\mathrm{thm}}$ layered by module; of six dependency edges, only one stays within its module, visualizing the low cohesion of~$0.107$ (\S\ref{sec:cohesion}). The \textcolor{blue!60!black}{blue} boundaries represent source files, in contrast to the \textcolor{green!40!black}{green} namespace boundaries of Figure~\ref{fig:ns-aggregation}; the two layers partially overlap (NMI~$\approx 0.71$; \S\ref{sec:ns-mod-cross}). All names have the \texttt{Nat.}\ prefix removed.}
\label{fig:module-region}
\end{figure}

\begin{definition}[Module cohesion]
\label{def:cohesion}
The \emph{cohesion} of a module~$m$ is
\[
  \mathrm{coh}(m) = \frac{|E_{\mathrm{int}}(m)|}{|E_{\mathrm{int}}(m)| + |E_{\mathrm{ext}}(m)|},
\]
with the convention that $\mathrm{coh}(m) = 0$ when $|E_{\mathrm{int}}(m)| + |E_{\mathrm{ext}}(m)| = 0$.
\end{definition}

\begin{table}[ht]
\centering
\caption{Global distribution of module cohesion across $6{,}993$ file-level modules.}
\label{tab:cohesion-global}
\begin{tabular}{lr}
\toprule
Statistic & Value \\
\midrule
Mean       & $0.107$ \\
Median     & $0.074$ \\
Std Dev    & $0.121$ \\
Max        & $1.000$ \\
Modules with $\mathrm{coh} = 0$ & $584$ / $6{,}993$ ($8.4\%$) \\
\bottomrule
\end{tabular}
\end{table}

Module cohesion is low throughout: the mean is~$0.107$ and the median is~$0.074$ (Table~\ref{tab:cohesion-global}). Internal edges constitute only about $10\%$ of a typical module's total connectivity. If module boundaries reflected logical structure, cohesion would cluster near~$1$; the observed value indicates that file boundaries and dependency structure are only weakly aligned.

\begin{remark}
\label{rem:zero-cohesion}
Among the $584$ modules ($8.4\%$) with zero cohesion, most are \emph{lemma collections}: declarations grouped by concept rather than by logical interdependence. Developers place ``all lemmas about~$X$'' in one file for discoverability, even when these lemmas depend on definitions across the library. Zero cohesion is not a defect but the structural signature of cognitive organization.
\end{remark}

\paragraph{Import edges vs.\ declaration-level dependencies.}
\label{sec:import_vs_declaration}

\noindent Aggregating $G_{\mathrm{thm}}$ to the module level yields the \emph{file-aggregated graph}~$G_{\mathrm{file}}$, in which an edge $m_1 \to m_2$ exists whenever some declaration in~$m_1$ cites one in~$m_2$ ($m_1 \neq m_2$). The gap is striking: $G_{\mathrm{file}}$ has $215{,}211$ edges versus $23{,}235$ in~$G_{\mathrm{module}}$---a $9.1\times$ amplification.

\begin{table}[ht]
\centering
\caption{Edge-level comparison of $G_{\mathrm{module}}$ (human-written imports) and $G_{\mathrm{file}}$ (declaration-aggregated dependencies).}
\label{tab:import-vs-file}
\begin{tabular}{lrr}
\toprule
Category & Edges & Percentage \\
\midrule
\multicolumn{3}{l}{\emph{$G_{\mathrm{module}}$ breakdown ($23{,}235$ edges):}} \\
\quad Active imports (in both $G_{\mathrm{module}}$ and $G_{\mathrm{file}}$) & $16{,}742$ & $72.1\%$ \\
\quad Unused imports ($G_{\mathrm{module}}$ only) & $6{,}493$ & $27.9\%$ \\
\midrule
\multicolumn{3}{l}{\emph{$G_{\mathrm{file}}$ breakdown ($215{,}211$ edges):}} \\
\quad Direct import exists & $16{,}742$ & $7.8\%$ \\
\quad Indirect, transitively reachable & $197{,}639$ & $91.8\%$ \\
\quad Indirect, not reachable & $830$ & $0.4\%$ \\
\bottomrule
\end{tabular}
\end{table}

From the import side, $72.1\%$ of import edges are \emph{active}---the importing file's declarations actually cite the imported module (Table~\ref{tab:import-vs-file}). The remaining $27.9\%$ provide transitive visibility only. From the declaration side, only $7.8\%$ of cross-file dependencies correspond to a direct import; the remaining $92.2\%$ are reached through transitive chains. The module graph functions as a sparse gateway: each file declares a handful of direct imports (mean out-degree $3.1$ in $G_{\mathrm{module}}$), but through transitivity these few edges open access to a vast dependency space (mean out-degree $30.8$ in $G_{\mathrm{file}}$).

\begin{remark}[Relationship to transitive redundancy]
\label{rem:unused-vs-redundant}
Section~\ref{sec:transitive-reduction} found $17.5\%$ of import edges transitively redundant; here, $27.9\%$ are unused. The gap ($27.9\% > 17.5\%$) has a precise interpretation: $78.4\%$ of unused imports are graph-theoretically essential---the sole path for transitive visibility---while $21.6\%$ are both unused and redundant.
\end{remark}

\paragraph{Import depth vs.\ usage depth.}

The preceding edge-count comparison treats all cross-module dependencies uniformly. We now refine the analysis by comparing the \emph{depth} of modules that a file directly imports in $G_{\mathrm{module}}$ with the depth of modules whose declarations it actually cites in $G_{\mathrm{thm}}$. For each module~$A$, let $\overline{d}_{\mathrm{imp}}(A)$ be the mean module depth of $A$'s direct imports and $\overline{d}_{\mathrm{use}}(A)$ the mean depth of the modules containing declarations cited by $A$'s own declarations. The \emph{depth difference} is $\Delta(A) = \overline{d}_{\mathrm{use}}(A) - \overline{d}_{\mathrm{imp}}(A)$.

Across the $6{,}861$ modules for which both quantities are defined, the mean depth difference is $\overline{\Delta} = -0.098$ (median $-0.069$, $\sigma = 0.445$). The distribution is roughly symmetric but slightly left-skewed: $54.2\%$ of modules have $\Delta < 0$ (usage shallower than imports), $40.3\%$ have $\Delta > 0$ (usage deeper), and $5.5\%$ have $\Delta = 0$. Actual usage thus reaches, on average, \emph{shallower} modules than those directly imported: through deep import chains, declarations access shallow algebraic and order-theoretic infrastructure.

\begin{table}[ht]
\centering
\caption{Mean depth difference $\overline{\Delta}$ by top-level directory (five most negative and five most positive among directories with ${\geq}30$ modules).}
\label{tab:depth-diff-by-dir}
\begin{tabular}{lrrrrr}
\toprule
Directory & $N$ & $\overline{d}_{\mathrm{imp}}$ & $\overline{d}_{\mathrm{use}}$ & $\overline{\Delta}$ & Median~$\Delta$ \\
\midrule
\multicolumn{6}{l}{\emph{Usage shallower than imports ($\Delta < 0$):}} \\
\module{CategoryTheory}      & 919 & 4.339 & 3.979 & $-0.360$ & $-0.333$ \\
\module{Geometry}             & 119 & 4.658 & 4.319 & $-0.339$ & $-0.364$ \\
\module{Algebra}              & 1210 & 4.487 & 4.309 & $-0.178$ & $-0.151$ \\
\module{Analysis}             & 702 & 4.522 & 4.358 & $-0.164$ & $-0.141$ \\
\module{MeasureTheory}        & 284 & 4.380 & 4.218 & $-0.162$ & $-0.141$ \\
\midrule
\multicolumn{6}{l}{\emph{Usage deeper than imports ($\Delta > 0$):}} \\
\module{LinearAlgebra}        & 324 & 4.054 & 4.234 & $+0.180$ & $+0.193$ \\
\module{NumberTheory}         & 207 & 4.233 & 4.375 & $+0.142$ & $+0.179$ \\
\module{RingTheory}           & 621 & 4.103 & 4.247 & $+0.145$ & $+0.177$ \\
\module{FieldTheory}          &  73 & 3.917 & 4.231 & $+0.314$ & $+0.307$ \\
\module{RepresentationTheory} &  32 & 3.949 & 4.262 & $+0.313$ & $+0.195$ \\
\bottomrule
\end{tabular}
\end{table}

Table~\ref{tab:depth-diff-by-dir} reveals a clear split. \module{CategoryTheory} ($\overline{\Delta} = -0.360$) stands out: its files import other category-theoretic modules at depth~${\sim}4.3$ but actually use declarations from modules at depth~${\sim}4.0$. This reflects the highly layered, self-contained nature of categorical formalization---deep categorical imports serve as portals to shallow definitional infrastructure. At the opposite extreme, \module{FieldTheory} ($\overline{\Delta} = +0.314$) imports modules at depth~${\sim}3.9$ but uses declarations at depth~${\sim}4.2$: its imports serve as cognitive shortcuts to deeper algebraic and ring-theoretic foundations. The positive-$\Delta$ directories (\module{RingTheory}, \module{NumberTheory}, \module{LinearAlgebra}, \module{FieldTheory}) are precisely those that build on deep algebraic infrastructure, confirming that the module graph functions as a sparse gateway to deeper dependency chains.

\subsection{Namespace--Declaration Interaction}
\label{sec:ns-decl-interaction}

The containment analysis of \S\ref{sec:containment-decay} measured how much declaration-level dependency traffic stays within namespace boundaries (see also Figure~\ref{fig:ns-aggregation}). We can now compare this with module cohesion to assess the relative constraining power of the two human organizational systems at the declaration level.

At depth~$1$, namespace containment is $22.2\%$: roughly one in five edges stays within the same depth-$1$ namespace. By depth~$2$, containment drops to $14.2\%$ and saturates near $12.6\%$ by depth~$3$ (Table~\ref{tab:containment-decay}). Module cohesion, the mean proportion of a module's edges that stay internal, is $0.107$---close to the file-level containment of $15.6\%$ reported in the same table.

The comparison reveals a consistent ordering: namespace containment at depth~$1$ ($22.2\%$) exceeds file-level containment ($15.6\%$), which in turn exceeds module cohesion ($10.7\%$). Namespaces, as conceptual categories, retain slightly more dependency traffic than files---a naming domain like \ns{Nat} captures more intra-group dependencies than the file \module{Data.Nat.Basic}, because the namespace spans multiple files and encloses a larger set of related declarations. But all three figures are far below~$50\%$: regardless of which human boundary one draws---naming hierarchy or file system---the majority of mathematical dependencies cross it.

\paragraph{Depth asymmetry.}
\label{sec:depth-asymmetry}

\noindent The $9.1\times$ edge amplification measures the \emph{quantitative} gap between $G_{\mathrm{module}}$ and declaration-level dependencies. A complementary analysis of naming depth reveals a \emph{directional} gap. For each module in $G_{\mathrm{module}}$, the \emph{module depth} is the number of dot-separated components (e.g., \module{Mathlib.Data.Nat.Basic} has depth~$4$). For each declaration in $G_{\mathrm{thm}}$, the \emph{declaration depth} is the number of namespace components (e.g., \decl{Nat.add\_comm} has depth~$1$). Table~\ref{tab:depth-asymmetry} compares the depth patterns of edges in the two graphs.

\begin{table}[ht]
\centering
\caption{Depth asymmetry: import edges vs.\ declaration-level dependencies.}
\label{tab:depth-asymmetry}
\begin{tabular}{lrr}
\toprule
& $G_{\mathrm{module}}$ & $G_{\mathrm{thm}}$ \\
\midrule
Same depth                         & $55.4\%$ & $51.4\%$ \\
Source deeper (deep $\to$ shallow) & $23.3\%$ & $37.6\%$ \\
Target deeper (shallow $\to$ deep) & $21.3\%$ & $11.0\%$ \\
Mean depth difference              & $+0.03$  & $+0.36$  \\
\bottomrule
\end{tabular}
\end{table}

Import edges are nearly symmetric: the mean depth difference is~$+0.03$, and the deep$\,\to\,$shallow and shallow$\,\to\,$deep proportions are balanced ($23.3\%$ vs.\ $21.3\%$). Developers import \emph{laterally}---from one depth-$4$ module to another---rather than vertically across naming layers (Remark~\ref{rem:import-depth-symmetry}). At the declaration level, the pattern inverts: the mean depth difference rises to~$+0.36$, and deep$\,\to\,$shallow edges ($37.6\%$) outnumber shallow$\,\to\,$deep edges ($11.0\%$) by more than $3{\times}$. Among declarations at depth~$\ge 2$, fully $69.3\%$ of outgoing edges target depth~$\le 1$---the shallow infrastructure layer of definitions like \decl{OfNat.ofNat}, \decl{Eq.refl}, and \decl{DFunLike.coe} that appeared as extreme hubs in \S\ref{sec:thm-centrality}.

The contrast is a structural consequence of transitivity. A single \texttt{import} brings an entire module into scope, providing transitive access to all upstream declarations. This compresses the vertical depth asymmetry into a lateral peer-to-peer pattern: importing \module{Data.Nat.Basic} implicitly conveys access to all the shallow infrastructure in \module{Init.Prelude}, without the importing module needing to reference depth-$0$ declarations directly. The module graph is thus not merely a sparse summary of declaration-level dependencies; it is a \emph{depth-compressed} summary that masks the gravitational pull of shallow infrastructure.

\subsection{Three Layers Combined}
\label{sec:three-layers}

Table~\ref{tab:three-layers} assembles the cross-level measurements into a single comparison.

\begin{table}[ht]
\centering
\caption{Cross-level alignment summary. Each row compares two organizational layers.}
\label{tab:three-layers}
\begin{tabular}{llrl}
\toprule
Comparison & Measure & Value & Interpretation \\
\midrule
Namespace vs.\ module            & NMI          & $0.71$  & High: human--human agreement \\
Namespace vs.\ Louvain community & NMI          & $0.34$  & Moderate: human $\neq$ logic \\
Namespace vs.\ declaration       & Containment  & $22.2\%$ & Low: naming $\neq$ dependency \\
Module vs.\ declaration          & Cohesion     & $0.107$ & Low: files $\neq$ logical units \\
\bottomrule
\end{tabular}
\end{table}

Three findings emerge.

\paragraph{Human coherence.} The NMI of~$0.71$ between namespaces and modules confirms that the two human organizational systems---naming and filing---largely agree. Developers who choose a file name and a namespace for a set of declarations are, in the majority of cases, making consistent decisions. The coherence is real but imperfect: $99.1\%$ of same-namespace declaration pairs reside in different files, because large namespaces like \ns{CategoryTheory} are partitioned across hundreds of compilation units.

\paragraph{Human--logic divergence.} Every measure of alignment between human organization and logical structure is low. Namespace vs.\ Louvain community: NMI~$= 0.34$. Namespace containment at depth~$1$: $22.2\%$. Module cohesion: $0.107$. The logical communities discovered by Louvain decomposition (\S\ref{sec:thm-community}) bear only moderate resemblance to the naming hierarchy, and both module boundaries and namespace boundaries fail to contain more than a fraction of the dependency traffic. The problem is not that humans organize the library poorly---their internal consistency (NMI~$0.71$) is high---but that mathematical logic does not respect human cognitive categories.

\paragraph{Depth compression.} The depth asymmetry analysis (\S\ref{sec:depth-asymmetry}) adds a directional dimension to this picture. At the declaration level, $37.6\%$ of edges flow from deeper to shallower namespaces (mean depth difference~$+0.36$), reflecting the gravitational pull of infrastructure definitions. At the module level, import edges are nearly symmetric (mean depth difference~$+0.03$). Transitivity compresses the vertical dependency structure into a lateral peer-to-peer pattern: $69.3\%$ of edges from deep declarations ($\mathrm{depth} \ge 2$) target shallow infrastructure ($\mathrm{depth} \le 1$), but the module graph hides this asymmetry behind a few peer-level \texttt{import} statements. The module graph is not merely sparse---it is \emph{depth-compressed}.

\paragraph{The central tension, quantified.} The cross-level analysis makes the central tension of this paper precise. The deterministic product---the dependency graph---records the logical structure of mathematics, which is inherently cross-cutting: a lemma about natural numbers may depend on ring theory, order theory, and decidability infrastructure. The human process---organizing declarations into files and namespaces---imposes boundaries that are cognitively useful but logically porous. The coherence of human decisions (NMI~$0.71$) and the incoherence of human--logic alignment (NMI~$0.34$, cohesion~$0.107$) together show that the divergence is systematic, not accidental: it arises from the fundamental mismatch between the hierarchical structure of human cognition and the cross-cutting structure of mathematical dependency.

\paragraph{The granularity equilibrium.}
The cross-level analysis implicitly rests on a design parameter that is itself a process-trace: the granularity of modules. \module{Mathlib}'s $6{,}993$ modules contain $308{,}129$ declarations, yielding a mean of approximately $44$ declarations per module. This ratio is not determined by mathematics---it is determined by the cognitive capacity of human developers to manage a coherent unit of code in a single file. Two thought experiments illuminate the point.

At one extreme, if each declaration occupied its own module, the module graph would coincide with the declaration graph: $G_{\mathrm{module}} \cong G_{\mathrm{thm}}$. Module cohesion would be trivially~$1.0$ (every module has no internal edges to lose), but the graph would have $308{,}129$ nodes and be unnavigable by humans. At the other extreme, if all declarations were placed in a single module, the module graph would collapse to a single vertex: no import structure, no compilation parallelism, no organizational signal. The current equilibrium---approximately $44$ declarations per module, cohesion~$0.107$, $9.1\times$ edge amplification---represents a compromise shaped by human cognitive constraints: the largest unit a developer can hold in working memory while writing and reviewing proofs. Figure~\ref{fig:granularity-spectrum} illustrates the three regimes.

\begin{figure}[ht]
\centering
\begin{subfigure}[b]{0.30\textwidth}
\centering
\begin{tikzpicture}[
  dot/.style={circle, fill=red!70!black, minimum size=3.5pt, inner sep=0pt},
  dep/.style={->, >=Stealth, red!70!black, thin},
  modbox/.style={draw=blue!50!black, fill=blue!6, rounded corners=2pt, semithick,
                 minimum size=0.6cm},
]
  \node[modbox] (m1) at (0,1.2) {};
  \node[modbox] (m2) at (1.1,1.2) {};
  \node[modbox] (m3) at (2.2,1.2) {};
  \node[modbox] (m4) at (0.55,0) {};
  \node[modbox] (m5) at (1.65,0) {};
  \node[dot] at (0,1.2) {};
  \node[dot] (d2) at (1.1,1.2) {};
  \node[dot] at (2.2,1.2) {};
  \node[dot] (d4) at (0.55,0) {};
  \node[dot] (d5) at (1.65,0) {};
  \draw[dep] (d4) -- (m1);
  \draw[dep] (d4) -- (d2);
  \draw[dep] (d5) -- (d2);
  \draw[dep] (d5) -- (m3);
  \draw[dep] (m1) -- (d2);
  \node[font=\scriptsize, text=gray] at (2.85,0.6) {$\cdots$};
\end{tikzpicture}
\caption{1 decl/module}
\label{fig:gran-fine}
\end{subfigure}%
\hfill
\begin{subfigure}[b]{0.36\textwidth}
\centering
\begin{tikzpicture}[
  dot/.style={circle, fill=red!70!black, minimum size=3.5pt, inner sep=0pt},
  dep/.style={->, >=Stealth, red!70!black, thin},
  depx/.style={->, >=Stealth, red!70!black, thin, dashed},
  modimp/.style={->, >=Stealth, blue!60!black, thick},
  modbox/.style={draw=blue!50!black, fill=blue!6, rounded corners=3pt, semithick},
]
  % Module A
  \node[modbox, minimum width=1.6cm, minimum height=1.4cm] at (0,0.6) {};
  \node[font=\tiny\sffamily, text=blue!60!black] at (0,1.55) {module $A$};
  \node[dot] (a1) at (-0.35,0.95) {};
  \node[dot] (a2) at (0.35,0.95) {};
  \node[dot] (a3) at (0,0.3) {};
  \draw[dep] (a3) -- (a1);
  \draw[dep] (a3) -- (a2);
  % Module B
  \node[modbox, minimum width=1.6cm, minimum height=1.0cm] at (2.8,0.6) {};
  \node[font=\tiny\sffamily, text=blue!60!black] at (2.8,1.35) {module $B$};
  \node[dot] (b1) at (2.45,0.6) {};
  \node[dot] (b2) at (3.15,0.6) {};
  \draw[dep] (b2) -- (b1);
  % Cross-module declaration edges
  \draw[depx] (a2) -- (b1);
  \draw[depx] (b2) to[out=-155,in=-25] (a3);
  % Blue import arrow
  \draw[modimp] (2.0,0.6) -- (0.8,0.6);
\end{tikzpicture}
\caption{${\sim}44$ decl/module (\module{Mathlib})}
\label{fig:gran-current}
\end{subfigure}%
\hfill
\begin{subfigure}[b]{0.30\textwidth}
\centering
\begin{tikzpicture}[
  dot/.style={circle, fill=red!70!black, minimum size=3.5pt, inner sep=0pt},
  dep/.style={->, >=Stealth, red!70!black, thin},
  modbox/.style={draw=blue!50!black, fill=blue!6, rounded corners=3pt, semithick},
]
  \node[modbox, minimum width=2.4cm, minimum height=1.6cm] at (1.1,0.6) {};
  \node[dot] (d1) at (0.35,1.0) {};
  \node[dot] (d2) at (1.1,1.0) {};
  \node[dot] (d3) at (1.85,1.0) {};
  \node[dot] (d4) at (0.7,0.25) {};
  \node[dot] (d5) at (1.5,0.25) {};
  \draw[dep] (d4) -- (d1);
  \draw[dep] (d4) -- (d2);
  \draw[dep] (d5) -- (d2);
  \draw[dep] (d5) -- (d3);
  \draw[dep] (d1) -- (d2);
\end{tikzpicture}
\caption{All decl/1 module}
\label{fig:gran-coarse}
\end{subfigure}
\caption{The granularity spectrum. \textcolor{blue!60!black}{Blue boxes} are modules; \textcolor{red!70!black}{red dots} are declarations; \textcolor{red!70!black}{solid red arrows} are intra-module dependencies; \textcolor{red!70!black}{dashed red arrows} are cross-module dependencies. (a)~One declaration per module: $G_{\mathrm{module}} \cong G_{\mathrm{thm}}$ ($308$K modules, unnavigable). (b)~Current \module{Mathlib}: multiple cross-module declaration edges collapse into a single \textcolor{blue!60!black}{blue import} arrow; internal edges are invisible at the module level. (c)~All declarations in one module: $G_{\mathrm{module}}$ collapses to a single vertex with no edges.}
\label{fig:granularity-spectrum}
\end{figure}

As AI systems become primary contributors to formal libraries, this equilibrium may shift. An AI author is not subject to the same working-memory constraint; it could operate productively with finer-grained modules optimized for compilation parallelism, or with coarser modules optimized for logical coherence. The direction of the shift---toward machine-optimal granularity rather than human-optimal granularity---would alter every cross-level metric reported in this section: cohesion, containment, edge amplification, and depth compression would all change, not because the mathematics changed, but because the production process changed. Tracking these metrics longitudinally as AI contributions grow will distinguish intrinsic mathematical structure from artifacts of the human production mode.

\paragraph{The namespace granularity spectrum.}
A parallel thought experiment applies to the \emph{namespace} axis. \module{Mathlib}'s $15{,}456$ leaf namespaces contain $308{,}129$ declarations, yielding a mean of approximately $20$ declarations per namespace. Again, this ratio is a product of human naming conventions, not of mathematical necessity.

At one extreme, if each declaration occupied its own namespace (flat naming: \ns{add\_comm}, \ns{dvd\_mul}, \ldots), containment at depth~$1$ would be trivially $0\%$---no two declarations could share a namespace boundary---and the naming hierarchy would carry no organizational information. At the other extreme, if all declarations shared a single namespace, containment would be trivially $100\%$ but the hierarchy would provide no finer-than-global grouping. The current Mathlib equilibrium---$6$ depth levels, $22.2\%$ containment at depth~$1$, rising to $86.9\%$ at depth~$3$---reflects a naming convention that balances logical grouping against navigational convenience. Figure~\ref{fig:ns-granularity-spectrum} illustrates these three regimes.

\begin{figure}[ht]
\centering
\begin{subfigure}[b]{0.30\textwidth}
\centering
\begin{tikzpicture}[
  dot/.style={circle, fill=red!70!black, minimum size=3.5pt, inner sep=0pt},
  dep/.style={->, >=Stealth, red!70!black, thin},
  nsbox/.style={draw=green!50!black, fill=green!6, rounded corners=2pt, semithick,
                minimum size=0.6cm},
]
  \node[nsbox] (n1) at (0,1.2) {};
  \node[nsbox] (n2) at (1.1,1.2) {};
  \node[nsbox] (n3) at (2.2,1.2) {};
  \node[nsbox] (n4) at (0.55,0) {};
  \node[nsbox] (n5) at (1.65,0) {};
  \node[dot] at (0,1.2) {};
  \node[dot] (d2) at (1.1,1.2) {};
  \node[dot] at (2.2,1.2) {};
  \node[dot] (d4) at (0.55,0) {};
  \node[dot] (d5) at (1.65,0) {};
  \draw[dep] (d4) -- (n1);
  \draw[dep] (d4) -- (d2);
  \draw[dep] (d5) -- (d2);
  \draw[dep] (d5) -- (n3);
  \draw[dep] (n1) -- (d2);
  \node[font=\scriptsize, text=gray] at (2.85,0.6) {$\cdots$};
\end{tikzpicture}
\caption{1 decl/namespace}
\label{fig:ns-gran-flat}
\end{subfigure}%
\hfill
\begin{subfigure}[b]{0.36\textwidth}
\centering
\begin{tikzpicture}[
  dot/.style={circle, fill=red!70!black, minimum size=3.5pt, inner sep=0pt},
  dep/.style={->, >=Stealth, red!70!black, thin},
  depx/.style={->, >=Stealth, red!70!black, thin, dashed},
  nsbox/.style={draw=green!50!black, fill=green!6, rounded corners=3pt, semithick},
  nsinr/.style={draw=green!50!black, fill=green!3, rounded corners=2pt, thin, dashed},
]
  % Outer namespace
  \node[nsbox, minimum width=3.2cm, minimum height=1.8cm] at (1.4,0.65) {};
  \node[font=\tiny\sffamily, text=green!50!black] at (1.4,1.8) {\ns{Nat}};
  % Inner namespace
  \node[nsinr, minimum width=1.3cm, minimum height=1.0cm] at (2.3,0.55) {};
  \node[font=\tiny\sffamily, text=green!50!black] at (2.3,1.25) {\ns{Prime}};
  % Declarations in Nat
  \node[dot] (a1) at (0.3,0.9) {};
  \node[dot] (a2) at (0.3,0.35) {};
  % Declarations in Nat.Prime
  \node[dot] (b1) at (1.9,0.55) {};
  \node[dot] (b2) at (2.7,0.55) {};
  % Intra-namespace edges
  \draw[dep] (a2) -- (a1);
  \draw[dep] (b2) -- (b1);
  % Cross-namespace edges
  \draw[depx] (b1) -- (a1);
  \draw[depx] (b2) to[out=170,in=-30] (a2);
\end{tikzpicture}
\caption{Hierarchical (\module{Mathlib})}
\label{fig:ns-gran-current}
\end{subfigure}%
\hfill
\begin{subfigure}[b]{0.30\textwidth}
\centering
\begin{tikzpicture}[
  dot/.style={circle, fill=red!70!black, minimum size=3.5pt, inner sep=0pt},
  dep/.style={->, >=Stealth, red!70!black, thin},
  nsbox/.style={draw=green!50!black, fill=green!6, rounded corners=3pt, semithick},
]
  \node[nsbox, minimum width=2.4cm, minimum height=1.6cm] at (1.1,0.6) {};
  \node[dot] (d1) at (0.35,1.0) {};
  \node[dot] (d2) at (1.1,1.0) {};
  \node[dot] (d3) at (1.85,1.0) {};
  \node[dot] (d4) at (0.7,0.25) {};
  \node[dot] (d5) at (1.5,0.25) {};
  \draw[dep] (d4) -- (d1);
  \draw[dep] (d4) -- (d2);
  \draw[dep] (d5) -- (d2);
  \draw[dep] (d5) -- (d3);
  \draw[dep] (d1) -- (d2);
\end{tikzpicture}
\caption{All decl/1 namespace}
\label{fig:ns-gran-single}
\end{subfigure}
\caption{The namespace granularity spectrum. \textcolor{green!50!black}{Green boxes} are namespaces; \textcolor{red!70!black}{red dots} are declarations; \textcolor{red!70!black}{solid red arrows} are intra-namespace dependencies; \textcolor{red!70!black}{dashed red arrows} cross namespace boundaries. (a)~One declaration per namespace: every edge crosses a boundary (containment~$= 0\%$). (b)~Current \module{Mathlib}: nested namespaces with partial containment. (c)~All declarations in one namespace: every edge is internal (containment~$= 100\%$, no organizational signal).}
\label{fig:ns-granularity-spectrum}
\end{figure}

Figures~\ref{fig:granularity-spectrum} and~\ref{fig:ns-granularity-spectrum} together reveal that library organization is a \emph{two-dimensional} design choice: module granularity (how many declarations per file) and namespace granularity (how many declarations per naming group) are independently adjustable parameters. \module{Mathlib} currently occupies one point in this design space---${\sim}44$ declarations per module, ${\sim}20$ per leaf namespace---but alternative design points are entirely feasible, and each would produce different cross-level statistics. The product (mathematical dependency) constrains neither axis; only the process (human or machine authorship) does.
